%% sections/j-counsel-routing.tex - H. R. 9510 Bill v4.0 (full bill) - Appendix J
%% Genre: legislative-counsel routing memo (TO/FROM/RE/DATE + amendatory-instruction
%% table + numbered form questions for HOLC). Source deliverable 10. Table 9.
%% Gray-scale Mermaid figure (TikZ): Figure 6 (the pre-introduction path by actor).
%% Single hyphens only; section symbol only in tables.
\billsec{APPENDIX J.}{LEGISLATIVE COUNSEL ROUTING MEMO (NON-OPERATIVE).}\phantomsection\label{toc:aJ}

\uscprov{1}{\textit{Independent research draft. Not an enacted law, not pending
legislation, and not legal advice; not endorsed by the FDA, HHS, the House Office of
the Legislative Counsel, or any Member of Congress. This memorandum is non-operative;
it transmits a draft for legislative-form review and carries no legal force of its
own.}}

\uscprov{0}{\textbf{TO:} House Office of the Legislative Counsel, via the sponsoring
office (Member of the House Committee on Energy and Commerce, Subcommittee on
Health)\\ \textbf{FROM:} CEO Kevin Kawchak, ChemicalQDevice\\ \textbf{RE:} Drafting
review and routing of H. R. 9510 (illustrative number; the Clerk assigns the real
number at introduction), 119th Congress, 2d Session, the Verification Before
Generation in Physical AI Oncology Trials Act of 2026, amending the Federal Food,
Drug, and Cosmetic Act (21 U.S.C. 301 et seq.), current through Public Law
119-93\\ \textbf{DATE:} June 7, 2026}

\uscprov{1}{\textbf{Purpose and routing.} This memo transmits a research draft of
H. R. 9510 for legislative-form review and asks the House Office of the Legislative
Counsel (HOLC) to confirm or correct the draft's amendatory structure before
introduction. HOLC provides nonpartisan, impartial, and confidential drafting
services to House Members and committees, not directly to the public, so this
proposal is submitted through the sponsoring Member's office, the proper channel for
engaging HOLC. We expect that several drafts may be required, and we ask that HOLC
treat the attached text as a starting draft to be revised through the office, not a
finished product \cite{house-holc, crs-r44001}.}

\uscprov{1}{The draft is in bill form and can become law after passage by both
chambers and presentment to the President; it is not a simple or concurrent
resolution. The designation H. R. 9510 is a placeholder; the Clerk assigns the
operative number only at introduction. The bill is paired with companion deliverables
in this package, including a section-by-section analysis, a Ramseyer-style
comparative print, legislative findings, a constitutional authority statement under
Rule XII, clause 7(c)(1), and a plain-English policy memo. Those documents are
supporting matter, not part of the bill text submitted for codification.}

\uscprov{1}{\textbf{Vehicle and positive-law question.} The bill amends the Federal
Food, Drug, and Cosmetic Act (FD\&C Act), 21 U.S.C. 301 et seq., current through
Public Law 119-93. The central form question is positive-law status. The FD\&C Act
has not been enacted into positive law as Title 21; Title 21 is a compilation of which
the FD\&C Act is the underlying and legally controlling text. The draft therefore
amends the underlying Act by the Act's own section numbers (for example, ``Section
515D of the Federal Food, Drug, and Cosmetic Act is amended''), giving the
corresponding Code citation parenthetically (for example, ``(21 U.S.C. 360e-5)'').
HOLC drafting guidance notes that positive-law status governs whether a measure
amends the Code provision or the underlying statute. We ask HOLC to confirm that this
Act-amends-Act approach, with parenthetical 21 U.S.C. citations, is correct and
consistently applied \cite{house-holc, olrc-usc21}.}

\uscprov{1}{\textbf{Amendatory structure.} Section 3 of the bill is the operative
provision. Subsection (a) inserts a new section 515D into subchapter V of the Act,
immediately after section 515C; subsection (b) makes the ten conforming amendments;
subsection (c) is the clerical amendment to the Act's table of contents; subsection
(d) is an Act-level rule of construction and savings provision; and subsection (e)
carries the effective date, the 365-day rulemaking deadline, and the 180-day
transition rule. Table 9 lists the new section and each conforming amendment, with its
FD\&C Act section, the parenthetical 21 U.S.C. citation, and the amendatory
instruction \cite{usc-360e4, pl-119-93}.}

{\footnotesize
\begin{xltabular}{\textwidth}{@{}L{2.2cm} L{1.6cm} L{1.6cm} Y@{}}
\hline\noalign{\vskip 0.04cm}
\textbf{Provision} & \textbf{FD\&C} & \textbf{21 U.S.C.} & \textbf{Amendatory instruction}\\
\hline\noalign{\vskip 0.04cm}
\endfirsthead
\hline\noalign{\vskip 0.04cm}
\textbf{Provision} & \textbf{FD\&C} & \textbf{21 U.S.C.} & \textbf{Amendatory instruction}\\
\hline\noalign{\vskip 0.04cm}
\endhead
\hline\noalign{\vskip 0.04cm}
\endlastfoot
New operative section & \S~515D & \S~360e-5 & Insert new section 515D in subchapter V
after section 515C, with subsections setting the verification, validation, and
uncertainty-quantification gate requirements\\
\hline\noalign{\vskip 0.04cm}
Conforming (1) & \S~201(h) & \S~321(h) & Confirm that robot-patient interaction code
and autonomous robot-control software are a device\\
\hline\noalign{\vskip 0.04cm}
Conforming (2) & \S~301 & \S~331 & Add a prohibited act for generation or execution in
violation of section 515D and for failure to document and attest\\
\hline\noalign{\vskip 0.04cm}
Conforming (3) & \S~501 & \S~351 & Deem a Physical AI device software function used
without a cleared section 515D record to be adulterated\\
\hline\noalign{\vskip 0.04cm}
Conforming (4) & \S~505F & \S~355g & Recognize section 515D records as real world
evidence for the program\\
\hline\noalign{\vskip 0.04cm}
Conforming (5) & \S~510(k) & \S~360(k) & Provide that a verification-governed change
under an established PCCP needs no new premarket notification\\
\hline\noalign{\vskip 0.04cm}
Conforming (6) & \S~513 & \S~360c & Graduate special controls by the device's degree
of autonomy\\
\hline\noalign{\vskip 0.04cm}
Conforming (7) & \S~515 & \S~360e & Require a PMA application to include the cleared
verification record and gate evidence\\
\hline\noalign{\vskip 0.04cm}
Conforming (8) & \S~515C & \S~360e-4 & Require a predetermined change control plan to
include an automated verification-before-change protocol\\
\hline\noalign{\vskip 0.04cm}
Conforming (9) & \S~520(o) & \S~360j(o) & Confirm such software is not within the
clinical-decision-support exclusion\\
\hline\noalign{\vskip 0.04cm}
Conforming (10) & \S~521 & \S~360k & Add a savings clause preserving State
human-in-the-loop and audit requirements\\
\hline\noalign{\vskip 0.04cm}
Clerical & Contents & (Act table) & Insert the ``Sec. 515D'' item after the item for
section 515C\\
\hline
\end{xltabular}}
\vspace{-0.8cm}
\figcaption{Table 9. The amendatory-instruction table: the new operative section,
the ten conforming amendments, and the clerical amendment, each with its FD\&C Act
section, its 21 U.S.C. citation, and the instruction.}

\uscprov{1}{We ask HOLC to confirm the amendatory verbs, the quotation form, and the
placement of inserted matter, in particular the ``is amended by inserting after
section 515C the following:'' construction and the quoted block for new section 515D,
and that each conforming amendment uses the correct point of insertion within its
existing subsection rather than a freestanding restatement.}

\uscprov{1}{\textbf{Cross-reference and disambiguation.} The bill turns on a
section-number distinction that is easy to misread, and we flag it for special
attention. Section 360(k) of Title 21 is the premarket notification (510(k)) pathway,
referenced in conforming amendment (5) and inside new section 515D and section 515C.
Section 360k of Title 21, with no parenthetical, is a different and freestanding
provision, State and local requirements respecting devices (preemption), amended by
conforming amendment (10) to add the savings clause. The two are not the same section
and must not be conflated. We have used ``section 360(k)'' only for the 510(k) pathway
and ``section 360k'' only for the preemption provision throughout the draft and the
comparative print. We ask HOLC to verify that every internal cross-reference,
including references from new section 515D to sections 515C, 515, 513, and 520(g), and
to the part 11, part 50, and part 312 frameworks, points to the intended target, and
that the conforming ``; and'' edits at the points of insertion match the current
text.}

\uscprov{1}{\textbf{Definitions and effective dates.} New section 515D defines the
operative terms, among them Physical AI system, robot agent, robot-patient
interaction code, generative artificial intelligence, AI-enabled device software
function, verification, validation, uncertainty quantification, verification
fraction, validation agreement, the coefficient-of-variation bound, the hard
catastrophe predicate, Physical AI Standard Level, Unification Standard Level, and
hash-chained audit trail. We ask HOLC to advise on reconciling these section-specific
definitions with the Act's general device definition in section 201(h) as amended, and
the new software-function terms with the clinical-decision-support exclusion in
section 520(o) as amended, so the local definitions neither contradict nor are
silently overridden by the general ones.}

\uscprov{1}{On timing, section 3(e) provides that the Act takes effect on the first
day of the first calendar quarter beginning not less than one year after enactment;
that the Secretary must issue final implementing regulations not later than 365 days
after enactment; and that an investigation already enrolling on the effective date has
180 days to come into compliance. The provision leaves the February 2, 2026, Quality
Management System Regulation effective date undisturbed. We ask HOLC to confirm that
the effective date, rulemaking deadline, and transition rule are mutually consistent
and free of gaps \cite{fr-qmsr-2024}.}

\uscprov{1}{\textbf{The pre-introduction path.} Figure 6 traces the route by actor.
The sponsoring office consolidates one authoritative text; experts and counsel
conduct human and legal review; the currency check in Appendix K re-pulls every
anchor; HOLC puts the measure in legislative form and attaches the constitutional
authority statement and the PAYGO note; the office recruits a sponsor and original
cosponsors; and the measure is introduced through the eHopper, after which the Clerk
assigns the number and refers it to Energy and Commerce.}

\begin{mermaidfig}
\node[mmstep,text width=40mm] (a) at (0,0) {Sponsoring office: consolidate one authoritative text};
\node[mmin,text width=40mm] (c) at (0,-1.4) {Experts and counsel: human and legal review};
\node[mmin,text width=40mm] (d) at (0,-2.8) {Currency check (Appendix K)};
\node[mmstep,text width=40mm] (e) at (0,-4.2) {HOLC: put in legislative form; attach CAS and PAYGO};
\node[mmstep,text width=40mm] (b) at (0,-5.6) {Recruit sponsor and original cosponsors};
\node[mmgoal,text width=40mm] (f) at (0,-7.0) {Introduce via the eHopper};
\node[mmlaw,text width=40mm] (g) at (0,-8.4) {Clerk assigns the number; refer to Energy and Commerce};
\draw[mmedge] (a)--(c); \draw[mmedge] (c)--(d); \draw[mmedge] (d)--(e); \draw[mmedge] (e)--(b); \draw[mmedge] (b)--(f); \draw[mmedge] (f)--(g);
\end{mermaidfig}
\vspace{-0.3cm}
\figcaption{Figure 6. The pre-introduction path by actor: consolidate, review, currency-check, route through the Office of the Legislative Counsel, recruit support, and introduce.}

\uscprov{1}{\textbf{Specific questions for Legislative Counsel.} We ask HOLC to
address the following form questions before the draft is suitable for introduction.}

\begin{enumerate}[leftmargin=2.0em,itemsep=2pt,topsep=2pt,label=(\arabic*)]
\item Should the bill amend the underlying FD\&C Act by its own section numbers with
parenthetical 21 U.S.C. citations, given that the Act is not positive-law Title 21,
and is that convention applied consistently throughout the draft?
\item Are the amendatory instructions and quotation form for new section 515D and for
each of the ten conforming amendments in correct legislative form?
\item Is the clerical amendment correctly framed to insert the ``Sec. 515D'' item
after the section 515C item in the Act's table of contents?
\item Are the section-number usages correct and unambiguous throughout, in particular
section 360(k) (the 510(k) premarket notification pathway) versus section 360k (State
and local preemption)?
\item Do the section 515D definitions need express reconciliation with section 201(h)
and with section 520(o), and if so, what coordinating language do you recommend?
\item Are the effective date, the 365-day rulemaking deadline, and the 180-day
transition rule internally consistent and properly drafted?
\item Should the savings clause in conforming amendment (10) and the rule of
construction in section 3(d) be coordinated to avoid any conflict at the preemption
boundary under section 360k?
\item Does the rule of construction correctly preserve the section 520(g)
investigational device exemption, the part 50 and part 312 frameworks, and the
section 515C authority, without inadvertently narrowing them?
\item Are there any drafting conventions, defined-term placements, or cross-reference
forms HOLC would change before this draft is suitable for introduction?
\end{enumerate}

\uscprov{1}{\textbf{Next steps.} We request HOLC's drafting review through the
sponsoring office and will incorporate any revisions into a new draft, preserving a
redline, version history, and a short explanation of any unresolved policy choices.
After HOLC's form review, the sponsoring office intends to seek procedural and
jurisdictional review from the House Parliamentarian's office on likely committee
referral and any points of order, confirm the sponsor and original cosponsors,
finalize the Rule XII constitutional authority statement, and proceed through the
current eHopper workflow at introduction. Only the final text approved by the
sponsor's office will be filed \cite{house-ehopper, crs-rs22477}.}
